\documentclass[12pt,notitlepage]{report}
% The shared notes preamble lives outside this repo.
\input{/Users/aryankulkarni/Documents/y3-notes/preamble.tex}

\title{Week 3: responses, explanatory variables and the GLM family}
\author{Aryan Kulkarni}
\date{MTH6134 group project}

\begin{document}
\maketitle

\paragraph{Findings.}
The data hold four kinds of response: a continuous one, a binary one, counts
and an ordered category. The first three map to the normal, Bernoulli and
Poisson members of the GLM family, and one algorithm fits all three. The
ordered response needs an extension (week 10).

\section{Which variables are responses}

A variable can explain a response only if its value is settled before the
response, or separately from it. Age, position and the previous three seasons
of injuries precede the 2023/24 season. Minutes played and injury are both
outcomes of that season, so neither is used to explain the other. Training
load is set during the season and is treated as explanatory, with the caution
that coaches may choose it in response to how a player is doing.

\begin{center}
  \small
  \begin{tabular}{llll}
    \toprule
    Response & Outcome & Distribution, link & Week \\
    \midrule
    \mintinline{r}|minutes_played| & continuous, largest value 3960 & normal, identity & 2, 10 \\
    \mintinline{r}|injured| & binary & Bernoulli, logit & 4, 5 \\
    \mintinline{r}|result == "Win"| & binary & Bernoulli, logit & 4 \\
    \mintinline{r}|goals_scored| & count & Poisson, log & 6 \\
    \mintinline{r}|goals_conceded| & count & Poisson, log & 6, 10 \\
    \mintinline{r}|prev_injuries| & count, observed 0 to 4 & Poisson, log & 6, 10 \\
    \mintinline{r}|result| & ordered, three levels & proportional odds & 10 \\
    \bottomrule
  \end{tabular}
\end{center}

The explanatory variables are \mintinline{r}|position|, \mintinline{r}|age|,
\mintinline{r}|training_load|, \mintinline{r}|prev_injuries| and
\mintinline{r}|bmi| for players, and \mintinline{r}|home_away| and
\mintinline{r}|opponent_rank| for matches. \mintinline{r}|team| and
\mintinline{r}|opponent| define groups (week 11). The categorical variables
are also cross-classified in weeks 8 and 9.

\section{What the counts look like}

A Poisson variable has variance equal to its mean, so the two are compared
for each count.

\begin{minted}{r}
source("../common.R")

counts <- list(
  goals_scored = matches$goals_scored,
  goals_conceded = matches$goals_conceded,
  prev_injuries = players$prev_injuries
)
moments <- sapply(counts, function(y) {
  c(mean = mean(y), variance = var(y), ratio = var(y) / mean(y), max = max(y))
})
out(round(moments, 2), "mean-variance")

count_data <- enframe(counts, name = "variable", value = "y") |> unnest(y)
observed <- count_data |>
  count(variable, y) |>
  group_by(variable) |>
  mutate(share = n / sum(n))
reference <- count_data |>
  group_by(variable) |>
  summarise(mu = mean(y)) |>
  expand_grid(y = 0:8) |>
  mutate(share = dpois(y, mu))

fig <- ggplot(observed, aes(x = y, y = share)) +
  geom_col(width = 0.6) +
  geom_point(data = reference, colour = orange, size = 1.6) +
  facet_wrap(~variable) +
  scale_x_continuous(breaks = 0:8) +
  labs(x = "Count", y = "Share of observations")
save_fig(fig, "counts", height = 2.4)
\end{minted}

\inputminted{text}{output/mean-variance.txt}

\begin{figure}[H]
  \centering
  \incfig{counts}
  \caption{Observed counts (bars) against a Poisson distribution with the
    same mean (points).}
  \label{fig:counts}
\end{figure}

Goals conceded sit close to the Poisson shape. Goals scored have a variance
1.5 times the mean, driven by one match with 8 goals; the Poisson assumption
concerns the distribution given the explanatory variables, so this ratio is
judged again in week 6 once home advantage and opponent strength are in the
model. Previous injuries have too many zeros and too many fours for a Poisson
variable, and the largest value is exactly 4. Week 10 returns to both
features.

\section{The three parts of a GLM}

A generalised linear model has
\begin{enumerate}
  \item a \emph{random part}: independent responses $Y_i$ from one
    exponential-family distribution with mean $\mu_i$,
  \item a \emph{systematic part}: a linear predictor
    $\eta_i = \bm{x}_i^\top \bm{\beta}$,
  \item a \emph{link}: a monotone function $g$ with $g(\mu_i) = \eta_i$.
\end{enumerate}
The exponential family has density
$f(y; \theta, \phi) = \exp\{(y\theta - b(\theta))/\phi + c(y, \phi)\}$, which
gives $\E(Y) = b'(\theta)$ and $\Var(Y) = \phi\, b''(\theta) = \phi V(\mu)$.
The variance is therefore a known function of the mean, and choosing a family
is choosing that function. The canonical link sets $\eta = \theta$.

\begin{center}
  \begin{tabular}{llll}
    \toprule
    Family & $V(\mu)$ & $\phi$ & Canonical link \\
    \midrule
    normal & 1 & $\sigma^2$ & identity, $\mu$ \\
    Bernoulli & $\mu(1 - \mu)$ & 1 & logit, $\log\{\mu/(1 - \mu)\}$ \\
    Poisson & $\mu$ & 1 & $\log \mu$ \\
    \bottomrule
  \end{tabular}
\end{center}

\paragraph{The linear model is a GLM.}
With the normal family and identity link, \mintinline{r}|glm| reproduces
\mintinline{r}|lm| exactly, and its dispersion estimate is $\hat\sigma^2$.

\begin{minted}{r}
fit_lm <- lm(minutes_played ~ position + training_load, data = players)
fit_glm <- glm(minutes_played ~ position + training_load,
               family = gaussian, data = players)
out(cbind(lm = coef(fit_lm), glm = coef(fit_glm)), "lm-is-glm")
out(c(sigma_squared = sigma(fit_lm)^2,
      dispersion = summary(fit_glm)$dispersion), "dispersion")
\end{minted}

\inputminted{text}{output/lm-is-glm.txt}
\inputminted{text}{output/dispersion.txt}

\section{How a GLM is fitted}

Maximum likelihood has no closed form outside the normal case. With a
canonical link the score equations are
$\bm{X}^\top(\bm{y} - \bm{\mu}) = \bm{0}$, and Newton's method for solving
them is a sequence of weighted least-squares fits. Given the current
$\bm{\beta}$, form the working response and weights
\[
  z_i = \eta_i + (y_i - \mu_i)\, g'(\mu_i),
  \qquad
  w_i = \frac{1}{V(\mu_i)\, g'(\mu_i)^2},
\]
regress $\bm{z}$ on $\bm{X}$ with weights $\bm{w}$, and repeat. For the
Poisson model with log link, $g'(\mu) = 1/\mu$, so $z = \eta + (y - \mu)/\mu$
and $w = \mu$. Eight iterations by hand match \mintinline{r}|glm|.

\begin{minted}{r}
X <- model.matrix(~ home_away + opponent_rank, data = matches)
y <- matches$goals_scored
beta <- c(log(mean(y)), 0, 0)
for (iteration in 1:8) {
  eta <- drop(X %*% beta)
  mu <- exp(eta)
  z <- eta + (y - mu) / mu
  beta <- drop(solve(t(X) %*% (mu * X), t(X) %*% (mu * z)))
}

fit_goals <- glm(goals_scored ~ home_away + opponent_rank,
                 family = poisson, data = matches)
out(cbind(by_hand = beta, glm = coef(fit_goals)), "irls")
\end{minted}

\inputminted{text}{output/irls.txt}

The final weighted fit also supplies the standard errors, since
$\widehat{\Var}(\hat{\bm{\beta}}) = \phi(\bm{X}^\top \bm{W} \bm{X})^{-1}$.

\section{Summary}

\begin{itemize}
  \item Injury is the binary response for week 4, goals and previous injuries
    are the counts for week 6, and the match result is an ordered response
    for week 10.
  \item A family is a statement about how the variance depends on the mean.
    The count plots already suggest where that statement will fail: previous
    injuries.
  \item Every model from here on is fitted by iterated weighted least
    squares, so the residuals, leverages and Cook's distances of week 2 all
    have GLM versions (week 5).
\end{itemize}

\end{document}
